\documentclass[12pt]{article}
\usepackage{amsmath}
\usepackage{amssymb}
\usepackage{geometry}
\geometry{letterpaper, margin=1in}

\title{Homework 2.3: Characterizations of Invertible Matrices}
\author{Tashfeen Omran}
\date{September 2026}

\begin{document}

\maketitle

\section{Part 1: Comprehensive Introduction, Context, and Prerequisites}

\subsection{Core Concepts}
The topic of ``Characterizations of Invertible Matrices'' centers entirely around a monumental result in linear algebra: the \textbf{Invertible Matrix Theorem (IMT)}. Up until this point, we have studied concepts like pivot positions, linear independence, the span of a set of vectors, null spaces (solutions to $A\mathbf{x} = \mathbf{0}$), and linear transformations mapping from $\mathbb{R}^n$ to $\mathbb{R}^n$. 

The core idea of this chapter is that for a square $n \times n$ matrix, \emph{all of these distinct concepts are inextricably linked}. If a matrix is "good" in one sense (e.g., its columns span the entire space), it is automatically "good" in all other senses (e.g., its columns are linearly independent, its determinant is non-zero, and it possesses an inverse). The IMT unifies the geometric perspective (transformations), the algebraic perspective (equations), and the structural perspective (matrix entries).

\subsection{Intuition and Derivation}
Why do all these properties imply one another? The intuition stems from the dimensions of the space. Imagine you are mapping a 3D space ($\mathbb{R}^3$) to another 3D space. 
If your transformation collapses any part of that space (e.g., squishing a 3D volume into a 2D plane), you lose information. Algebraically, this means $A\mathbf{x} = \mathbf{0}$ has non-trivial solutions (multiple vectors get crushed to the origin). Because you are restricted to three dimensions, squishing the space implies you cannot possibly reach every point in the target 3D space—thus, the columns cannot span $\mathbb{R}^3$. 

Conversely, if you preserve the dimension (linear independence, $A\mathbf{x} = \mathbf{0}$ only having the trivial solution), you have enough "directions" to reach every point in the target space (spanning). Because no information is lost and every point is reachable, you can perfectly reverse the process. This reverse mapping is exactly what the inverse matrix, $A^{-1}$, represents. 

\subsection{Historical Context and Motivation}
Historically, mathematicians in the 18th and 19th centuries were deeply concerned with solving massive systems of linear equations—often arising from physics, astronomy, and later, economics. Solving these systems by hand was painstakingly slow. Mathematicians needed to know \emph{before} doing the calculation whether a unique solution even existed. If a system had infinitely many solutions or no solutions, attempting to find "the" unique inverse was a waste of time. The Invertible Matrix Theorem evolved as a checklist: if you could easily verify just \emph{one} property (like checking for linearly dependent columns), you could definitively conclude the system's solvability without doing the grueling arithmetic of Gaussian elimination.

\subsection{Key Formulas and Theorems}
\begin{itemize}
    \item \textbf{Definition of an Inverse:} For an $n \times n$ matrix $A$, if there exists an $n \times n$ matrix $C$ such that $AC = CA = I_n$ (where $I_n$ is the identity matrix), then $C = A^{-1}$.
    \item \textbf{Inverse of a $2 \times 2$ Matrix:} For $A = \begin{bmatrix} a & b \\ c & d \end{bmatrix}$, if $ad - bc \neq 0$, then 
    \[
    A^{-1} = \frac{1}{ad - bc} \begin{bmatrix} d & -b \\ -c & a \end{bmatrix}
    \]
    \item \textbf{The Invertible Matrix Theorem (IMT):} Let $A$ be a square $n \times n$ matrix. The following are all true or all false:
    \begin{enumerate}
        \item[(a)] $A$ is an invertible matrix.
        \item[(c)] $A$ has $n$ pivot positions.
        \item[(d)] $A\mathbf{x} = \mathbf{0}$ has only the trivial solution ($\mathbf{x} = \mathbf{0}$).
        \item[(e)] The columns of $A$ form a linearly independent set.
        \item[(g)] The equation $A\mathbf{x} = \mathbf{b}$ has at least one solution for each $\mathbf{b} \in \mathbb{R}^n$.
        \item[(h)] The columns of $A$ span $\mathbb{R}^n$.
        \item[(i)] The linear transformation $\mathbf{x} \mapsto A\mathbf{x}$ maps $\mathbb{R}^n$ onto $\mathbb{R}^n$.
        \item[(f)] The linear transformation $\mathbf{x} \mapsto A\mathbf{x}$ is one-to-one.
    \end{enumerate}
    \textit{(Note: This is a condensed list of the most utilized statements from the theorem).}
\end{itemize}

\subsection{Prerequisites}
To master this topic, you must have a rock-solid understanding of:
\begin{enumerate}
    \item \textbf{Row Reduction (Gaussian Elimination):} Ability to find pivot positions.
    \item \textbf{Linear Independence and Span:} Understanding what it means for vectors to be multiples of each other, and what it means to fill a space.
    \item \textbf{Matrix Multiplication and Algebra:} Knowing that $AB \neq BA$ in general, but associativity $A(BC) = (AB)C$ holds.
\end{enumerate}


\section{Part 2: Detailed Homework Solutions}

\textit{Note: As requested, a representative selection of problems covering every distinct type introduced in the chapter has been solved to ensure comprehensive coverage without redundancy.}

\subsection{Type 1: Visual Invertibility (Exercises 2 \& 4)}

\textbf{Problem 2:} Determine if the matrix is invertible. Use as few calculations as possible.
\[
A = \begin{bmatrix} -4 & 6 \\ 6 & -9 \end{bmatrix}
\]

\textbf{Solution:}
We inspect the columns of matrix $A$. Let $\mathbf{c}_1 = \begin{bmatrix} -4 \\ 6 \end{bmatrix}$ and $\mathbf{c}_2 = \begin{bmatrix} 6 \\ -9 \end{bmatrix}$.
Notice that column 2 is a scalar multiple of column 1. Specifically:
\[
-1.5 \cdot \mathbf{c}_1 = -1.5 \cdot \begin{bmatrix} -4 \\ 6 \end{bmatrix} = \begin{bmatrix} 6 \\ -9 \end{bmatrix} = \mathbf{c}_2
\]
Because one column is a scalar multiple of another, the columns form a linearly dependent set. According to the Invertible Matrix Theorem (statement e), if the columns of an $n \times n$ matrix are not linearly independent, the matrix is not invertible. 
\textbf{Final Answer:} The matrix is \textbf{not invertible}.

\vspace{0.5cm}

\textbf{Problem 4:} Determine if the matrix is invertible. Use as few calculations as possible.
\[
A = \begin{bmatrix} -7 & 0 & 4 \\ 3 & 0 & -1 \\ 2 & 0 & 9 \end{bmatrix}
\]

\textbf{Solution:}
We inspect the columns of the matrix. The second column is a zero vector, $\mathbf{c}_2 = \begin{bmatrix} 0 \\ 0 \\ 0 \end{bmatrix}$. 
Any set of vectors containing the zero vector is automatically linearly dependent. By the Invertible Matrix Theorem (statement e), since the columns are linearly dependent, the matrix is not invertible. (Alternatively, a matrix with a column of zeros can never have a pivot in that column, meaning it cannot have $n$ pivots, violating statement c).
\textbf{Final Answer:} The matrix is \textbf{not invertible}.

\subsection{Type 2: IMT Logical Implications (Exercises 11 \& 16)}

\textbf{Problem 11:} (True/False) If the equation $A\mathbf{x} = \mathbf{0}$ has only the trivial solution, then $A$ is row equivalent to the $n \times n$ identity matrix. Justify your answer. (Assume $A$ is an $n \times n$ matrix).

\textbf{Solution:}
The statement ``$A\mathbf{x} = \mathbf{0}$ has only the trivial solution'' is statement (d) of the Invertible Matrix Theorem. 
The statement ``$A$ is row equivalent to the $n \times n$ identity matrix'' is statement (b) of the IMT.
Because $A$ is a square matrix, the IMT applies. The theorem dictates that if any one statement is true, all other statements are true. Therefore, (d) implies (b).
\textbf{Final Answer:} \textbf{True}. By the Invertible Matrix Theorem, statement (d) implies statement (b).

\vspace{0.5cm}

\textbf{Problem 16:} (True/False) If the equation $A\mathbf{x} = \mathbf{b}$ has at least one solution for each $\mathbf{b}$ in $\mathbb{R}^n$, then the solution is unique for each $\mathbf{b}$. (Assume $A$ is an $n \times n$ matrix).

\textbf{Solution:}
The premise ``$A\mathbf{x} = \mathbf{b}$ has at least one solution for each $\mathbf{b}$ in $\mathbb{R}^n$'' is statement (g) of the IMT. 
Because statement (g) is true, $A$ is an invertible matrix.
If $A$ is invertible, we can multiply both sides of $A\mathbf{x} = \mathbf{b}$ by $A^{-1}$:
\begin{align*}
    A^{-1}(A\mathbf{x}) &= A^{-1}\mathbf{b} \\
    (A^{-1}A)\mathbf{x} &= A^{-1}\mathbf{b} \\
    I_n\mathbf{x} &= A^{-1}\mathbf{b} \\
    \mathbf{x} &= A^{-1}\mathbf{b}
\end{align*}
Since matrix multiplication yields a single, unique resulting vector, $A^{-1}\mathbf{b}$ is unique. Therefore, the solution exists and is unique. (Alternatively, (g) implies (f), which means the transformation is one-to-one, guaranteeing uniqueness).
\textbf{Final Answer:} \textbf{True}.

\subsection{Type 3: Theoretical Matrix Properties (Exercise 23)}

\textbf{Problem 23:} Can a square matrix with two identical columns be invertible? Why or why not?

\textbf{Solution:}
Let $A$ be an $n \times n$ matrix. Assume column $i$ and column $j$ are identical, so $\mathbf{c}_i = \mathbf{c}_j$ (for $i \neq j$).
We can form a linear combination of the columns of $A$ that equals the zero vector without all weights being zero. Specifically, we can set the weight for $\mathbf{c}_i$ to $1$, the weight for $\mathbf{c}_j$ to $-1$, and all other weights to $0$:
\[
1 \cdot \mathbf{c}_i + (-1) \cdot \mathbf{c}_j + \sum_{k \neq i,j} 0 \cdot \mathbf{c}_k = \mathbf{c}_i - \mathbf{c}_i = \mathbf{0}
\]
Because there is a non-trivial solution to $A\mathbf{x} = \mathbf{0}$ (specifically, a vector $\mathbf{x}$ with $1$ in the $i$-th entry, $-1$ in the $j$-th entry, and $0$ elsewhere), the columns are linearly dependent. By the IMT, $A$ cannot be invertible.
\textbf{Final Answer:} \textbf{No.} The identical columns cause the columns to be linearly dependent, which by the IMT means the matrix is not invertible.

\subsection{Type 4: Matrix Algebra Proofs (Exercise 35)}

\textbf{Problem 35:} Show that if $AB$ is invertible, so is $A$. You cannot use Theorem 6(b), because you cannot assume that $A$ and $B$ are invertible. [Hint: There is a matrix $W$ such that $ABW = I$. Why?] (Assume $A$ and $B$ are $n \times n$ matrices).

\textbf{Solution:}
Let $A$ and $B$ be $n \times n$ matrices. We are given that the product matrix $(AB)$ is invertible.
By the definition of invertibility, there exists an $n \times n$ inverse matrix. Let us call this inverse matrix $W$. Therefore:
\[
(AB)W = I
\]
Because matrix multiplication is associative, we can regroup the terms on the left side:
\[
A(BW) = I
\]
Let $C$ be an $n \times n$ matrix defined as $C = BW$. Substituting $C$ into our equation gives:
\[
AC = I
\]
According to the Invertible Matrix Theorem (specifically statement k implying statement a), if there exists an $n \times n$ matrix $C$ such that $AC = I$, then $A$ is an invertible matrix.
\textbf{Final Answer:} The proof is complete. By establishing a right inverse $C = BW$, the IMT guarantees $A$ is fully invertible.

\subsection{Type 5: Invertible Linear Transformations (Exercise 41)}

\textbf{Problem 41:} $T$ is a linear transformation from $\mathbb{R}^2$ into $\mathbb{R}^2$. Show that $T$ is invertible and find a formula for $T^{-1}$.
\[
T(x_1, x_2) = (-9x_1 + 7x_2, 4x_1 - 3x_2)
\]

\textbf{Solution:}
First, we find the standard matrix $A$ for the linear transformation $T$. The transformation of the standard basis vectors is:
\begin{align*}
    T(1, 0) &= (-9(1) + 7(0), 4(1) - 3(0)) = \begin{bmatrix} -9 \\ 4 \end{bmatrix} \\
    T(0, 1) &= (-9(0) + 7(1), 4(0) - 3(1)) = \begin{bmatrix} 7 \\ -3 \end{bmatrix}
\end{align*}
Thus, the standard matrix $A$ is:
\[
A = \begin{bmatrix} -9 & 7 \\ 4 & -3 \end{bmatrix}
\]
To show $T$ is invertible, we must show $A$ is invertible by checking its determinant:
\[
\det(A) = (-9)(-3) - (7)(4) = 27 - 28 = -1
\]
Since the determinant is non-zero ($-1 \neq 0$), matrix $A$ is invertible, meaning $T$ is an invertible linear transformation.
Next, we calculate $A^{-1}$ using the $2 \times 2$ formula:
\[
A^{-1} = \frac{1}{-1} \begin{bmatrix} -3 & -7 \\ -4 & -9 \end{bmatrix} = \begin{bmatrix} 3 & 7 \\ 4 & 9 \end{bmatrix}
\]
The formula for $T^{-1}$ is governed by $A^{-1}$. Let $\mathbf{x} = (x_1, x_2)$. 
\[
T^{-1}(\mathbf{x}) = A^{-1} \begin{bmatrix} x_1 \\ x_2 \end{bmatrix} = \begin{bmatrix} 3 & 7 \\ 4 & 9 \end{bmatrix} \begin{bmatrix} x_1 \\ x_2 \end{bmatrix} = \begin{bmatrix} 3x_1 + 7x_2 \\ 4x_1 + 9x_2 \end{bmatrix}
\]
\textbf{Final Answer:} $T$ is invertible because its standard matrix has a determinant of $-1 \neq 0$. The inverse transformation is $T^{-1}(x_1, x_2) = (3x_1 + 7x_2, 4x_1 + 9x_2)$.


\section{Part 3: In-Depth Analysis of Problems and Techniques}

\subsection{A) Problem Types and General Approach}
\begin{enumerate}
    \item \textbf{Setup Only / Visual Inspection Problems (Exercises 1-10):}
    These problems test your ability to spot structural flaws that violate the IMT. The strategy is to scan for obvious linear dependence: columns of zeros, rows of zeros, or columns that are clear scalar multiples of each other. No heavy row reduction is needed.
    
    \item \textbf{Theoretical Implications / Logic Chains (Exercises 11-20):}
    These are True/False questions that require mapping a textual phrase (e.g., "onto") to its IMT letter (e.g., statement i), and then confirming that the premise and conclusion are both on the IMT list. The general strategy: Check if the matrix is square. If it is square, all IMT properties imply one another.
    
    \item \textbf{Matrix Algebra Proofs (Exercises 33-40):}
    These problems require manipulating symbolic matrices without knowing their entries. The core strategy relies on associativity $A(BC) = (AB)C$ and the definition of the inverse. You generally try to algebraically isolate the matrix in question alongside another matrix to form the identity matrix $I$, thereby proving invertibility via definition.

    \item \textbf{Linear Transformation Inverses (Exercises 41-48):}
    The approach here is algorithmic: extract the standard matrix $A$ from the equations, verify $\det(A) \neq 0$, compute $A^{-1}$, and then translate $A^{-1}$ back into a system of equations for the coordinates.
\end{enumerate}

\subsection{B) Key Algebraic and Matrix Manipulations}
\begin{itemize}
    \item \textbf{Exploiting Linear Dependence:} In Problem 2, recognizing that $\mathbf{c}_2 = -1.5 \mathbf{c}_1$ immediately solved the problem. Identifying scalar multiples prevents unnecessary Gaussian elimination.
    \item \textbf{Associative Regrouping Trick:} In Problem 35, the trick was regrouping $(AB)W$ to $A(BW)$. Matrix multiplication is not commutative ($AB \neq BA$), so you cannot swap the order, but you \emph{can} change where the parentheses go. This is the most crucial algebraic trick in matrix proofs.
    \item \textbf{The $2 \times 2$ Inverse Formula:} In Problem 41, instead of setting up an augmented matrix $[A \mid I]$ and row reducing, utilizing the $1/(ad-bc)$ trick is a massive time-saver for $2 \times 2$ matrices.
\end{itemize}


\section{Part 4: "Cheatsheet" and Tips for Success}

\subsection{Important Formulas and Theorems}
\begin{itemize}
    \item \textbf{The $2 \times 2$ Inverse Check:} For $\begin{bmatrix} a & b \\ c & d \end{bmatrix}$, invertibility exists $\iff ad-bc \neq 0$.
    \item \textbf{Reversal Law of Inverses:} $(AB)^{-1} = B^{-1}A^{-1}$. (Note the order flips!)
    \item \textbf{Transpose Inverse Law:} $(A^T)^{-1} = (A^{-1})^T$.
\end{itemize}

\subsection{Tips, Tricks, and Shortcuts}
\begin{itemize}
    \item \textbf{The "Zero" Cheat Code:} If a square matrix has a row of all zeros or a column of all zeros, it is \textbf{never} invertible. You don't need to calculate anything else.
    \item \textbf{The Triangular Matrix Shortcut:} For an upper or lower triangular matrix, the matrix is invertible if and only if all entries on the main diagonal are non-zero. The determinant is simply the product of the diagonal.
\end{itemize}

\subsection{Common Pitfalls and Traps}
\begin{itemize}
    \item \textbf{The Non-Square Trap:} The Invertible Matrix Theorem \textbf{ONLY} applies to square ($n \times n$) matrices. A $3 \times 2$ matrix might have linearly independent columns, but it cannot be invertible, nor can its columns span $\mathbb{R}^3$. Do not apply the IMT to rectangular matrices!
    \item \textbf{The Commutative Trap:} When taking the inverse of an algebraic equation like $AX = B$, if you multiply by $A^{-1}$, you must do it on the \emph{same side} on both sides. $A^{-1}AX = A^{-1}B \implies X = A^{-1}B$. Writing $X = BA^{-1}$ is a fatal error.
    \item \textbf{Zero Divisor Fallacy:} If $AB = \mathbf{0}$, it does \emph{not} mean $A = \mathbf{0}$ or $B = \mathbf{0}$. However, if $AB = \mathbf{0}$ and $A$ is \emph{invertible}, then $B$ must be $\mathbf{0}$ (proved by multiplying both sides by $A^{-1}$).
\end{itemize}


\section{Part 5: Conceptual Synthesis and The "Big Picture"}

\subsection{A) Thematic Connections}
The central theme of this topic is \textbf{Equivalence}. In earlier chapters, solving systems of equations, analyzing vectors spanning a space, and examining geometric transformations felt like three separate fields of math. The IMT synthesizes them. It shows that "algebraic uniqueness" (only the trivial solution exists), "geometric completeness" (spanning the whole space), and "structural integrity" ($n$ pivots) are exactly the same concept viewed through different lenses. 

\subsection{B) Forward and Backward Links}
\begin{itemize}
    \item \textbf{Backward Link:} This chapter relies heavily on row reduction (Chapter 1). A matrix fails the IMT if row reduction creates a row of zeros (meaning fewer than $n$ pivots).
    \item \textbf{Forward Link:} The IMT is the bedrock for finding \textbf{Eigenvalues and Eigenvectors} later in the course. An eigenvalue $\lambda$ exists only if the matrix $(A - \lambda I)$ is \emph{not} invertible. Thus, you will spend half of your future curriculum deliberately looking for matrices that violate the IMT in order to find their eigenvalues. It also perfectly sets the stage for Determinants (Chapter 3), where invertibility is summarized into a single number.
\end{itemize}


\section{Part 6: Real-World Application and Modeling}

\subsection{A) Concrete Scenarios in Finance and Economics}
\begin{enumerate}
    \item \textbf{Markowitz Mean-Variance Portfolio Optimization:} Quants at hedge funds seek to construct portfolios that maximize return for a given level of risk. This requires calculating the inverse of the covariance matrix of asset returns ($\Sigma$). If two assets are perfectly correlated (e.g., one stock is exactly tracking another), the covariance matrix has linearly dependent columns. By the IMT, the matrix becomes non-invertible (singular), and the optimization algorithm will crash because a unique optimal portfolio weight cannot be calculated.
    
    \item \textbf{Ordinary Least Squares (OLS) in Econometrics:} When building multi-factor models for securities pricing (like the Fama-French 3-factor model), economists use OLS regression. The formula for the beta coefficients requires the calculation of $(X^T X)^{-1}$. If the independent variables are perfectly collinear (e.g., including both 'total revenue' and 'total sales' which are perfectly identical columns), $X^T X$ violates the IMT, has no inverse, and the regression model mathematically breaks down.
    
    \item \textbf{Implicit Finite Difference Methods for Option Pricing:} When pricing exotic derivatives, financial engineers numerically solve the Black-Scholes Partial Differential Equation on a grid. In implicit schemes, stepping backward in time requires solving a massive linear system $A\mathbf{x} = \mathbf{b}$ at every time step. The matrix $A$ (usually a tridiagonal matrix representing transition probabilities) MUST be invertible. If it isn't, the option price cannot be resolved.
\end{enumerate}

\subsection{B) Model Problem Setup: Portfolio Optimization}
\textbf{Scenario:} We want to find the optimal weights for a portfolio of $n$ assets to minimize risk for a target expected return.
\textbf{Variables:}
\begin{itemize}
    \item $\mathbf{w}$: An $n \times 1$ vector of portfolio weights (the unknowns we want to find).
    \item $\boldsymbol{\mu}$: An $n \times 1$ vector of expected returns for the assets.
    \item $\Sigma$: The $n \times n$ covariance matrix of asset returns.
\end{itemize}
\textbf{The Model:} 
To find the minimum variance portfolio, one solves the Lagrangian formulation which eventually requires solving a linear system. A simplified version of calculating the global minimum variance portfolio weights requires solving:
\[
\Sigma \mathbf{w} = \mathbf{1}
\]
where $\mathbf{1}$ is a vector of ones. 
To find the weights $\mathbf{w}$, we must calculate:
\[
\mathbf{w} = \Sigma^{-1} \mathbf{1}
\]
\textbf{The IMT Application:} The solution $\mathbf{w}$ can only be uniquely determined if $\Sigma$ is invertible. If an analyst accidentally includes linearly dependent assets, $\Sigma$ will fail the IMT, $\Sigma^{-1}$ will not exist, and the portfolio solver will fail.


\section{Part 7: Common Variations and Untested Concepts}

\subsection{Missing Concept: Ill-Conditioned Matrices and Numerical Stability}
While the provided homework covers absolute invertibility (true or false), real-world systems deal with matrices that are \emph{almost} not invertible, known as \textbf{ill-conditioned matrices}. This concept was referenced in the textbook's "Numerical Notes" and Exercises 49-53, but not explicitly tested in our selected theoretical problems.

An ill-conditioned matrix is mathematically invertible, but practically problematic. A tiny change in the input (due to roundoff error or slight data variations) causes massive, chaotic changes in the output vector $\mathbf{x}$. The \emph{Condition Number} measures this sensitivity.

\textbf{Example Problem:}
Consider the system $A\mathbf{x} = \mathbf{b}$:
\[
\begin{bmatrix} 1 & 1 \\ 1 & 1.001 \end{bmatrix} \begin{bmatrix} x_1 \\ x_2 \end{bmatrix} = \begin{bmatrix} 2 \\ 2 \end{bmatrix}
\]
\textit{Solution:}
Solving this yields $x_1 = 2$, $x_2 = 0$. 

Now, introduce a tiny measurement error of $0.001$ in the second equation's constant:
\[
\begin{bmatrix} 1 & 1 \\ 1 & 1.001 \end{bmatrix} \begin{bmatrix} x_1 \\ x_2 \end{bmatrix} = \begin{bmatrix} 2 \\ 2.001 \end{bmatrix}
\]
\textit{Solution:}
Subtracting row 1 from row 2 gives $0.001x_2 = 0.001 \implies x_2 = 1$. Then $x_1 = 1$. 
\textbf{Takeaway:} A microscopic change of $0.001$ in the right-hand side caused the solution to jump from $(2,0)$ to $(1,1)$. Because the columns of $A$ are extremely close to being linearly dependent, $A$ is ill-conditioned. In financial modeling, trusting the inverse of an ill-conditioned matrix can lead to catastrophic trading errors.


\section{Part 8: Advanced Diagnostic Testing: "Find the Flaw"}

Below are five solutions to relevant matrix problems. Each contains a subtle, fatal flaw. Your task as the grader is to identify the error, explain why it is wrong in one sentence, and provide the correct solution.

\subsection{Problem 1: Invertibility Analysis}
\textbf{Question:} Let $A = \begin{bmatrix} 1 & 0 \\ 0 & 1 \\ 0 & 0 \end{bmatrix}$. Does the linear transformation $\mathbf{x} \mapsto A\mathbf{x}$ map $\mathbb{R}^2$ onto $\mathbb{R}^3$? Determine this by deciding if $A$ is invertible using the IMT.

\textbf{Flawed Solution:} 
We check the columns of $A$. The columns $\begin{bmatrix} 1 \\ 0 \\ 0 \end{bmatrix}$ and $\begin{bmatrix} 0 \\ 1 \\ 0 \end{bmatrix}$ are clearly linearly independent. Since the columns of $A$ are linearly independent, statement (e) of the Invertible Matrix Theorem is true. Therefore, $A$ is invertible. Because $A$ is invertible, statement (i) is true, meaning the transformation maps $\mathbb{R}^2$ onto $\mathbb{R}^3$. 

\textbf{Grader Correction:}
\begin{itemize}
    \item \textbf{The Flaw:} The student applied the Invertible Matrix Theorem to a $3 \times 2$ matrix.
    \item \textbf{Explanation:} The Invertible Matrix Theorem is strictly reserved for square ($n \times n$) matrices; rectangular matrices can never be invertible.
    \item \textbf{Correct Solution:} Since $A$ is $3 \times 2$, it is not square and cannot be invertible. A transformation mapping $\mathbb{R}^2$ to $\mathbb{R}^3$ can at most span a 2D plane in $\mathbb{R}^3$. It is impossible to map onto $\mathbb{R}^3$.
\end{itemize}

\subsection{Problem 2: Matrix Algebra Expansion}
\textbf{Question:} Simplify the expression $(AB)^{-1} A$, assuming $A$ and $B$ are invertible $n \times n$ matrices.

\textbf{Flawed Solution:}
First, we distribute the inverse over the product. 
$(AB)^{-1} A = (A^{-1} B^{-1}) A$.
Because matrix multiplication is associative, we rearrange to group the $A$ terms:
$A^{-1} (B^{-1} A)$.
Since matrices don't always commute, we can't easily cancel $A^{-1}$ and $A$. Therefore, the simplified form is $A^{-1} B^{-1} A$.

\textbf{Grader Correction:}
\begin{itemize}
    \item \textbf{The Flaw:} The student wrote $(AB)^{-1} = A^{-1}B^{-1}$.
    \item \textbf{Explanation:} The inverse of a product is the product of the inverses in reverse order: $(AB)^{-1} = B^{-1}A^{-1}$.
    \item \textbf{Correct Solution:} $(AB)^{-1} A = (B^{-1} A^{-1}) A = B^{-1} (A^{-1} A) = B^{-1} I = B^{-1}$.
\end{itemize}

\subsection{Problem 3: Solving a Matrix Equation}
\textbf{Question:} Solve for the $n \times n$ matrix $X$ in the equation $AX - B = C$. Assume $A$ is invertible.

\textbf{Flawed Solution:}
First, add $B$ to both sides:
$AX = C + B$.
To isolate $X$, we divide both sides by the matrix $A$:
$X = \frac{C + B}{A}$.

\textbf{Grader Correction:}
\begin{itemize}
    \item \textbf{The Flaw:} The student divided by a matrix.
    \item \textbf{Explanation:} Matrix division is entirely undefined; you must multiply by the inverse matrix, paying strict attention to left vs. right multiplication.
    \item \textbf{Correct Solution:} $AX = C + B$. Multiply both sides \emph{on the left} by $A^{-1}$: $A^{-1}(AX) = A^{-1}(C + B) \implies X = A^{-1}(C + B)$.
\end{itemize}

\subsection{Problem 4: Determinants and Solutions}
\textbf{Question:} Let $A$ be a $2 \times 2$ matrix where $\det(A) = 0$. Determine how many solutions the equation $A\mathbf{x} = \mathbf{b}$ has for a specific vector $\mathbf{b} = \begin{bmatrix} 1 \\ 1 \end{bmatrix}$.

\textbf{Flawed Solution:}
The determinant of $A$ is zero. According to the Invertible Matrix Theorem, if the determinant is zero, the matrix is not invertible. If a matrix is not invertible, then the system of equations it represents has broken down. Therefore, the equation $A\mathbf{x} = \mathbf{b}$ has absolutely no solutions. 

\textbf{Grader Correction:}
\begin{itemize}
    \item \textbf{The Flaw:} The student assumed that a non-invertible matrix always results in a system with zero solutions.
    \item \textbf{Explanation:} If $\det(A) = 0$, the system $A\mathbf{x} = \mathbf{b}$ either has no solutions OR infinitely many solutions, depending entirely on whether $\mathbf{b}$ is in the column space of $A$.
    \item \textbf{Correct Solution:} Since $A$ is singular, the columns of $A$ do not span $\mathbb{R}^2$, but rather form a line (or a point). If $\mathbf{b} = \begin{bmatrix} 1 \\ 1 \end{bmatrix}$ lies on that line, there are infinitely many solutions. If it does not, there are no solutions. We need the actual entries of $A$ to decide between 0 or $\infty$ solutions.
\end{itemize}

\subsection{Problem 5: Sum of Inverses}
\textbf{Question:} Suppose $A$ and $B$ are invertible $n \times n$ matrices. Prove that $(A + B)$ is invertible and its inverse is $A^{-1} + B^{-1}$.

\textbf{Flawed Solution:}
To prove a matrix is an inverse, we multiply the matrix by its proposed inverse and check if it equals the identity matrix $I$.
We test the proposed inverse $(A^{-1} + B^{-1})$ against $(A + B)$:
$(A + B)(A^{-1} + B^{-1}) = A(A^{-1}) + B(B^{-1}) = I + I = 2I$.
Since it yields an identity-like matrix, $(A+B)$ is invertible, we just have to adjust the scalar. The inverse is $\frac{1}{2}(A^{-1} + B^{-1})$.

\textbf{Grader Correction:}
\begin{itemize}
    \item \textbf{The Flaw:} The student improperly distributed the FOIL expansion of matrix sums.
    \item \textbf{Explanation:} $(A+B)(A^{-1}+B^{-1}) = AA^{-1} + AB^{-1} + BA^{-1} + BB^{-1} = I + AB^{-1} + BA^{-1} + I$, which generally does not simplify to $I$.
    \item \textbf{Correct Solution:} In general, $(A + B)$ is \textbf{not} necessarily invertible just because $A$ and $B$ are. For example, if $B = -A$, then $A+B = \mathbf{0}$, which has no inverse. Furthermore, there is no simple formula to split an inverse over a sum; $(A+B)^{-1} \neq A^{-1} + B^{-1}$. The premise of the problem itself is false.
\end{itemize}

\end{document}